El tubo de difracción de electrones se conectó a una fuente de potencial, la 
cual se varió entre 3000 y 5000 \si{\V}, con el fin de acelerar los electrones 
y dirigirlos hacia una lámina de grafito que actúa como rejilla de difracción.

Una vez difractados, los electrones generan un patrón de interferencia que 
puede observarse en una pantalla. En el caso del tubo cilíndrico, esta pantalla 
se encuentra a una distancia de 13.5 \si{\centi\meter}, como se muestra en la 
\cref{fig:Interference-pattern-cylinder}.

Para este experimento, también se utilizó un segundo montaje con un tubo de 
forma esférica, en el cual la pantalla está situada a 65 \si{\mili\meter}, 
distancia que corresponde a la longitud de la curva, como se ilustra en la 
\cref{fig:Interference-pattern-sphere}.

\begin{figure}[htbp!]
  \begin{center}
    \includegraphics[width=\linewidth, page=2]{../Images/fig:Interference-pattern-cylinder.jpg}
    \caption{Patrón de interferencia en el tubo con forma cilíndrica.}
    \label{fig:Interference-pattern-cylinder}
  \end{center}
\end{figure}

\begin{figure}[htbp!]
  \begin{center}
    \includegraphics[width=\linewidth, page=2]{../Images/fig:Interference-pattern-sphere.jpg}
  	\caption{Patrón de interferencia en el tubo con forma esférica.}
    \label{fig:Interference-pattern-sphere}
  \end{center}
\end{figure} 

Una vez obtenido el patrón de interferencia, se utilizó un calibrador de 
Vernier para medir los diámetros de los dos anillos que se generan. Este paso 
se debe hacer para cada cambio del potencial de la fuente, estos cambios fueron
de 500 \si{\V}.
